\documentclass[12pt]{article}

\usepackage[utf8]{inputenc}
\usepackage{geometry}
\geometry{a4paper, margin=1in}
\linespread{1.5}
\usepackage{amsmath}
\usepackage{amsfonts}
\usepackage{enumitem}
\usepackage{hyperref}
\hypersetup{hidelinks, colorlinks=true, allcolors=black}

\title{Lecture 1: The Solow Model\\
\large 3143 Textbook Chapters 4--6}
\author{Harlly Zhou}
\date{\today}

\begin{document}

\maketitle

\noindent\textbf{Reference:} Garín, Lester, and Sims, \textit{Intermediate Macroeconomics}, Chapters 4--6.

\tableofcontents
\newpage

%=============================================================================
\section{Chapter 4: Facts About Economic Growth}
%=============================================================================

\subsection{Overview}
Before modeling growth, we describe stylized facts. Kaldor (1957) listed six facts characterizing economic growth across developed countries. These remain a reasonable description of the data.

\subsection{The Kaldor Facts}
\begin{enumerate}
    \item \textbf{Output per worker} grows at a sustained, roughly constant rate over long periods.
    \begin{itemize}
        \item US: average growth rate $\approx 1.7\%$ per year.
        \item \textbf{Rule of 70}: Years to double $\approx 70 / g$ where $g$ is the annual growth rate (\%).
        \item At 1.7\%, GDP per worker doubles about every 41 years.
    \end{itemize}
    
    \item \textbf{Capital per worker} grows at a sustained, approximately constant rate.
    \item \textbf{Capital-to-output ratio} $K/Y$ is roughly constant over long periods.
    \item \textbf{Labor's share of income} is roughly constant:
    \[
        LABSH_t = \frac{w_t N_t}{Y_t}, \qquad CAPSH_t = 1 - LABSH_t
    \]
    US: labor share $\approx 0.65$, capital share $\approx 0.35$.
    \item \textbf{Rate of return on capital} $R_t$ is relatively constant. From $CAPSH_t = R_t K_t / Y_t$:
    \[
        R_t = (1 - LABSH_t) \frac{Y_t}{K_t}
    \]
    \item \textbf{Real wages} grow at a sustained, approximately constant rate (similar to output per worker).
\end{enumerate}

\subsection{Summary of Kaldor Facts}
Wages, output per worker, and capital per worker grow at approximately the same sustained rate; the return on capital is approximately constant. \textbf{Implication:} Economic growth benefits labor (real wages rise) and not capital (return on capital is roughly constant).

\subsection{Cross-Country Facts}
\begin{enumerate}
    \item \textbf{Enormous variation in income across countries.} GDP per capita (PPP) varies by orders of magnitude (e.g., US $\approx$ 36$\times$ Mali).
    \item \textbf{Growth miracles and growth disasters.} East Asian tigers (South Korea, Taiwan, China, Botswana) grew spectacularly; some African countries saw income decline.
    \item \textbf{Strong positive correlation} between income per capita and human capital (e.g., years of education).
\end{enumerate}

%=============================================================================
\section{Chapter 5: The Basic Solow Model}
%=============================================================================

\subsection{Overview}
The Solow model (Solow 1956) is the principal framework for understanding long-run growth and cross-country income differences. It assumes a constant saving rate and does not explicitly model microeconomic optimization.

\subsection{Production}
\begin{itemize}
    \item Aggregate production: $Y_t = A F(K_t, N_t)$ where $A$ is productivity (exogenous).
    \item Properties: $F_K > 0$, $F_N > 0$; $F_{KK} < 0$, $F_{NN} < 0$; $F_{KN} > 0$; constant returns to scale.
    \item Cobb-Douglas: $F(K_t, N_t) = K_t^{\alpha} N_t^{1-\alpha}$, $0 < \alpha < 1$.
    \item Firm FOCs: $w_t = A F_N(K_t, N_t)$, $R_t = A F_K(K_t, N_t)$.
\end{itemize}

\subsection{Consumption, Investment, and Capital Accumulation}
\begin{itemize}
    \item Resource constraint: $Y_t = C_t + I_t$.
    \item Constant saving rate: $I_t = s Y_t$, $C_t = (1-s) Y_t$, $0 < s < 1$.
    \item Capital accumulation: $K_{t+1} = I_t + (1-\delta) K_t$ where $\delta \in (0,1)$ is the depreciation rate.
    \item Labor supplied inelastically; $N_t$ exogenous and constant.
\end{itemize}

\subsection{Central Equation (Per Worker)}
Define $k_t = K_t / N_t$, $y_t = Y_t / N_t$, and $f(k_t) = F(k_t, 1)$. The law of motion for capital per worker:
\begin{equation}
    k_{t+1} = s A f(k_t) + (1-\delta) k_t
\end{equation}
For Cobb-Douglas: $f(k_t) = k_t^{\alpha}$, so
\begin{equation}
    k_{t+1} = s A k_t^{\alpha} + (1-\delta) k_t
\end{equation}
Other variables: $y_t = A k_t^{\alpha}$, $c_t = (1-s) A k_t^{\alpha}$, $i_t = s A k_t^{\alpha}$, $R_t = A f'(k_t)$, $w_t = A[f(k_t) - k_t f'(k_t)]$.

\subsection{Graphical Analysis}
\begin{itemize}
    \item Plot $k_{t+1}$ vs.\ $k_t$. The curve starts at origin, is increasing and concave (diminishing returns).
    \item \textbf{Inada conditions}: $\lim_{k \to 0} f'(k) = \infty$, $\lim_{k \to \infty} f'(k) = 0$.
    \item \textbf{Steady state} $k^*$: where $k_{t+1} = k_t$ (intersection with 45-degree line).
    \item \textbf{Convergence}: From any $k_0 > 0$, $k_t \to k^*$ over time.
\end{itemize}

\subsection{Alternative Depiction}
$\Delta k_{t+1} = k_{t+1} - k_t = s A f(k_t) - \delta k_t$. Steady state where investment $sAf(k_t)$ equals depreciation $\delta k_t$.

\subsection{Algebraic Steady State (Cobb-Douglas)}
\begin{equation}
    k^* = \left( \frac{sA}{\delta} \right)^{\frac{1}{1-\alpha}}, \qquad
    y^* = A (k^*)^{\alpha}, \qquad
    c^* = (1-s) y^*, \qquad
    i^* = s y^*
\end{equation}
$k^*$ is increasing in $s$ and $A$, decreasing in $\delta$.

\subsection{Experiments}
\begin{enumerate}
    \item \textbf{Increase in $s$}: Curve shifts up; new steady state $k_1^* > k_0^*$. Capital grows over time toward new steady state. Output growth jumps up temporarily, then returns to zero. Consumption initially falls; investment jumps up.
    \item \textbf{Increase in $A$}: Similar to increase in $s$. Higher productivity $\Rightarrow$ higher steady state $k^*$ and $y^*$.
\end{enumerate}

\subsection{Key Implication}
\textbf{Output will not forever grow faster} if the saving rate increases. There is a temporary burst of growth, but the economy returns to a steady state with zero growth in output per worker. Sustained growth in the basic Solow model requires something else (e.g., productivity growth).

\subsection{The Golden Rule}
The \textbf{golden rule} saving rate $s^{GR}$ maximizes steady-state consumption per worker. From $c^* = (1-s) A (k^*)^{\alpha}$ and $s A (k^*)^{\alpha} = \delta k^*$, the golden rule satisfies $f'(k^{GR}) = \delta$ (marginal product of capital equals depreciation). For Cobb-Douglas, $s^{GR} = \alpha$.

%=============================================================================
\section{Chapter 6: The Augmented Solow Model}
%=============================================================================

\subsection{Introducing Productivity and Population Growth}
\begin{itemize}
    \item Productivity grows at rate $g$: $A_{t+1} = (1+g) A_t$.
    \item Population (labor) grows at rate $n$: $N_{t+1} = (1+n) N_t$.
    \item Define \textbf{effective labor}: $A_t N_t$ grows at rate $g + n$ (approximately).
    \item Define \textbf{capital per effective worker}: $\kappa_t = K_t / (A_t N_t)$.
    \item Output per effective worker: $\phi_t = Y_t / (A_t N_t) = \kappa_t^{\alpha}$.
\end{itemize}

\subsection{Law of Motion for $\kappa_t$}
\begin{equation}
    \kappa_{t+1} = \frac{s \kappa_t^{\alpha} + (1-\delta) \kappa_t}{(1+g)(1+n)} \approx s \kappa_t^{\alpha} - (g + n + \delta) \kappa_t \quad \text{(in continuous time)}
\end{equation}
In steady state: $s \kappa^{\alpha} = (g + n + \delta) \kappa$, so
\begin{equation}
    \kappa^* = \left( \frac{s}{g + n + \delta} \right)^{\frac{1}{1-\alpha}}
\end{equation}

\subsection{Output and Consumption Per Capita}
\begin{itemize}
    \item Output per capita: $y_t = A_t \kappa_t^{\alpha}$. In steady state, $\kappa_t = \kappa^*$ constant, so $y_t$ grows at rate $g$ (productivity growth).
    \item \textbf{Key result:} Sustained growth in output per capita comes from productivity growth $g$, not from capital accumulation. Capital deepening ($\kappa$) is constant in steady state.
\end{itemize}

\subsection{Experiments}
\begin{enumerate}
    \item \textbf{Increase in $s$}: Higher $\kappa^*$; temporarily higher growth; long-run growth rate of $y_t$ unchanged (still $g$).
    \item \textbf{Increase in $A$} (or $g$): Shifts curve; higher $\kappa^*$; permanently higher growth rate of output per capita.
\end{enumerate}

\subsection{The Golden Rule (Augmented Model)}
The golden rule $\kappa^{GR}$ maximizes steady-state consumption per effective worker. Condition: $f'(\kappa^{GR}) = g + n + \delta$. For Cobb-Douglas, $s^{GR} = \alpha$ (same as basic model).

\subsection{Will Economic Growth Continue Indefinitely?}
Yes, as long as productivity $A_t$ continues to grow. The augmented Solow model predicts that output per capita grows at rate $g$ in the long run. The source of sustained growth is \textbf{productivity growth}, not capital accumulation.

%=============================================================================
\section{Summary}
%=============================================================================
\begin{itemize}
    \item \textbf{Ch.\ 4}: Kaldor facts---output, capital, wages grow; $K/Y$, labor share, return on capital roughly constant. Cross-country: large income variation, growth miracles/disasters, human capital--income correlation.
    \item \textbf{Ch.\ 5}: Basic Solow---constant $s$, $Y = AF(K,N)$, $K_{t+1} = I_t + (1-\delta)K_t$. Steady state $k^*$; convergence. Higher $s$ $\Rightarrow$ higher $k^*$ but only \emph{temporary} growth. Golden rule: $s^{GR} = \alpha$.
    \item \textbf{Ch.\ 6}: Augmented Solow---exogenous $g$ (productivity) and $n$ (population). Capital per effective worker $\kappa$ constant in steady state. \textbf{Sustained growth in $y$ comes from productivity growth $g$}, not capital accumulation. Policy implication: focus on productivity, not just saving.
\end{itemize}

\end{document}
